\documentclass[11pt]{article}
 
% ============================================================
% Packages
% ============================================================
\usepackage[a4paper,margin=1in]{geometry}
\usepackage[T1]{fontenc}
\usepackage[utf8]{inputenc}
% \usepackage{lmodern}      % uncomment if available (recommended)
 
\usepackage{amsmath, amssymb, amsthm}
\usepackage{mathtools}
\usepackage{graphicx}
\usepackage{xcolor}
\usepackage{enumitem}
\usepackage{booktabs}
\usepackage{longtable}
\usepackage{hyperref}
\usepackage{listings}
% \usepackage{microtype}    % uncomment if available (recommended)
\usepackage{setspace}
\usepackage{titlesec}
\usepackage{parskip}
\usepackage{array}
\usepackage{multirow}
\usepackage{mdframed}
 
% ============================================================
% Airier document layout
% ============================================================
\linespread{1.08}
\setlength{\parskip}{0.75em}
\setlength{\parindent}{0pt}
\setlength{\emergencystretch}{2em}
 
\titlespacing*{\section}{0pt}{3.0ex plus 1ex minus .2ex}{1.6ex plus .2ex}
\titlespacing*{\subsection}{0pt}{2.4ex plus .8ex minus .2ex}{1.1ex plus .2ex}
\titlespacing*{\subsubsection}{0pt}{2.0ex plus .6ex minus .2ex}{0.9ex plus .2ex}
 
\setlist[itemize]{topsep=0.6em,itemsep=0.35em,parsep=0.15em,leftmargin=2.2em}
\setlist[enumerate]{topsep=0.6em,itemsep=0.35em,parsep=0.15em,leftmargin=2.2em}
\setlist[description]{topsep=0.6em,itemsep=0.4em,parsep=0.2em,leftmargin=2.8em,style=nextline}
 
\renewcommand{\arraystretch}{1.3}
\setlength{\tabcolsep}{7pt}
 
% ============================================================
% Listings configuration
% ============================================================
\lstset{
    basicstyle=\ttfamily\small,
    breaklines=true,
    frame=single,
    columns=fullflexible,
    keepspaces=true,
    showstringspaces=false,
    tabsize=4,
    captionpos=b,
    aboveskip=1.0em,
    belowskip=1.0em,
    xleftmargin=1.0em,
    xrightmargin=1.0em
}
 
% ============================================================
% Theorem-like environments
% ============================================================
\newtheorem{remark}{Remark}[section]
\newtheorem{definition}{Definition}[section]
\newtheorem{contract}{Contract}
 
% Framed box for pipeline stages
\newmdenv[
  linewidth=0.8pt,
  innerleftmargin=10pt,
  innerrightmargin=10pt,
  innertopmargin=8pt,
  innerbottommargin=8pt,
  backgroundcolor=gray!6
]{stagebox}
 
% ============================================================
% Macros
% ============================================================
\newcommand{\RIGHT}{\texttt{RIGHT}}
\newcommand{\WRONG}{\texttt{WRONG}}
\newcommand{\dinst}{d_t^{\mathrm{inst}}}
\newcommand{\dbest}{d_t^{\mathrm{best}}}
\newcommand{\rtarget}{r_t^{\mathrm{target}}}
\newcommand{\rN}{r_N^{\mathrm{target}}}
\newcommand{\Xtest}{\mathcal{X}_{\mathrm{test}}}
 
% ============================================================
% Title
% ============================================================
\title{%
  \textbf{GPBO Separation of Concerns --- Experiment A}\\[0.4em]
  \large Reference: Methodology, Data Pipeline, and Code Architecture
}
\author{Benjamin Wilhelm}
\date{May 2026}
 
\begin{document}
 
\maketitle
 
\begin{abstract}
This document is the primary reference for Experiment A of the
\emph{Gaussian-Process Bayesian Optimization Separation of Concerns} project.
It covers the scientific motivation and experimental protocol (\S\S\,\ref{sec:intro}--\ref{sec:protocol}),
the precise definitions of all collected quantities and metrics (\S\,\ref{sec:quantities}),
the three-stage data pipeline from raw trajectories to aggregate statistics (\S\,\ref{sec:pipeline}),
the calibration diagnostic procedure (\S\,\ref{sec:calibration}),
the statistical analysis plan (\S\,\ref{sec:stats}),
a step-by-step notebook analysis roadmap (\S\,\ref{sec:notebook}),
the intended figures (\S\,\ref{sec:figures}),
and the code architecture (\S\,\ref{sec:code}).
Experimental contracts are collected in \S\,\ref{sec:contracts}.
Planned extensions are noted explicitly throughout; no unimplemented feature is described
as existing.
\end{abstract}
 
\tableofcontents
\clearpage
 
% ============================================================
\section{Introduction and Scientific Objective}
\label{sec:intro}
% ============================================================
 
\subsection{The Separation of Concerns Principle in BO}
 
Bayesian optimization rests on a clean conceptual separation:
\begin{itemize}
  \item the \emph{surrogate model} (typically a GP) encodes beliefs about the
        unknown function $f$;
  \item the \emph{acquisition function} encodes the decision rule that translates
        those beliefs into a query location.
\end{itemize}
The \textbf{Separation of Concerns (SoC)} principle states that the surrogate should
model the raw unknown phenomenon, while task-specific structure --- targets, priors,
constraints, costs, preferences --- should enter only through the acquisition function.
 
Violating SoC means baking task-specific transformations into the GP training targets.
This is a widespread practice (e.g., training a GP on $|f(x)-T|$ to handle
target-matching, or on a scalarization for multi-objective problems) and its empirical
consequences are largely unstudied.
 
\subsection{Research Questions}
 
The overarching project aims to characterize, across multiple SoC scenarios, the
empirical cost of the violation: does it degrade model quality? calibration? optimization
performance? or all three in a causally connected chain?
 
\medskip
\textbf{Experiment A} is the first and cleanest instantiation, focusing on the
\emph{target-matching} problem
\[
  \min_{x \in \mathcal{X}} |f(x) - T|.
\]
Three concrete questions are tested:
\begin{enumerate}[label=(Q\arabic*)]
  \item Does training the GP on $|f(x)-T|$ degrade optimization performance?
  \item Does it produce worse-calibrated uncertainty estimates?
  \item Can we visually and statistically identify the \emph{mechanism} of failure?
\end{enumerate}
 
The expected diagnostic chain is:
\[
  \text{SoC violation}
  \;\Rightarrow\;
  \text{model miscalibration near target}
  \;\Rightarrow\;
  \text{confidence-bound failure}
  \;\Rightarrow\;
  \text{acquisition failure}
  \;\Rightarrow\;
  \text{optimization failure.}
\]
 
% ============================================================
\section{Experiment A: Target Matching}
\label{sec:expA}
% ============================================================
 
\subsection{Problem Definition}
 
Let $\mathcal{X} \subset \mathbb{R}^d$ be a compact domain and $f:\mathcal{X}\to\mathbb{R}$
an unknown black-box function.
We observe noisy evaluations
\[
  y_t = f(x_t) + \varepsilon_t, \qquad \varepsilon_t \sim \mathcal{N}(0,\sigma^2),
\]
independently across iterations.
The goal is to find $x$ such that $f(x)$ is close to a prescribed target $T \in \mathbb{R}$.
 
\subsection{Why Target Matching Exposes the SoC Issue}
 
Suppose $f$ is smooth and well-modeled by a GP with a smooth kernel.
Define the transformed objective $g(x) = |f(x)-T|$.
Even though $f$ is smooth, $g$ develops \emph{kinks} along the target level set
$\{x: f(x) = T\}$: the absolute-value function is non-differentiable at zero.
Training a GP directly on $g$ therefore imposes a smooth covariance structure on a
non-smooth function. The SoC-compliant alternative avoids this by keeping the GP on $f$
and placing $T$ only in the acquisition.
 
\subsection{Benchmark Functions}
 
All benchmarks inherit from \texttt{SyntheticBenchmark}; see
\S\,\ref{sec:code-benchmarks} for the class hierarchy.
 
The registry \texttt{BENCHMARK\_REGISTRY} (in \texttt{benchmarks/\_\_init\_\_.py})
exposes eight benchmarks, grouped into the four ``standard'' analytic functions
and four minimal/diagnostic 1D target-matching benchmarks whose transformed
objective $|f-T|$ has deliberately controlled geometry (a single cusp or
several misleading near-target valleys).

\begin{center}
\begin{tabular}{lllll}
\toprule
Key & Name & $d$ & Domain (default) & Default $\sigma$ \\
\midrule
\texttt{sine\_trend\_1d} & Sine-trend 1D & 1 & $[0,1]$ & $0.05$ \\
\texttt{forrester\_1d}  & Forrester 2008 & 1 & $[0,1]$ & $0.5$ \\
\texttt{branin\_2d}     & Branin--Hoo    & 2 & $[-5,10]\times[0,15]$ & $2.0$ \\
\texttt{hartmann\_6d}   & Hartmann 6D    & 6 & $[0,1]^6$ & $0.05$ \\
\midrule
\texttt{linear\_1d}              & Linear monotone        & 1 & $[0,1]$ & $0.0$ \\
\texttt{exponential\_1d}         & Normalised exponential & 1 & $[0,1]$ & $0.0$ \\
\texttt{gaussian\_trap\_1d}      & Gaussian near-target trap & 1 & $[0,1]$ & $0.0$ \\
\texttt{damped\_oscillatory\_trap\_1d} & Damped oscillatory trap & 1 & $[0,1]$ & $0.0$ \\
\bottomrule
\end{tabular}
\end{center}

The first four are smooth/multimodal references. \texttt{linear\_1d}
($f(x)=x$) and \texttt{exponential\_1d}
($f(x)=(e^{ax}-1)/(e^a-1)$) are bijective, so every attainable target has a
unique cusp in $|f-T|$. \texttt{gaussian\_trap\_1d}
($f(x)=x-A\,e^{-(x-c)^2/2s^2}$) and \texttt{damped\_oscillatory\_trap\_1d}
($f(x)=x+a\,e^{-\lambda x}\sin(2\pi m x)$) add \emph{local} near-target valleys
to the transformed objective, stress-testing the \WRONG{} surrogate.

\textbf{Default $\sigma$} is a module-level per-benchmark constant
(\texttt{DEFAULT\_NOISE\_STD\_*}); it is always overridable by the YAML
\texttt{noise\_std} field. The default noiseless runs in \texttt{configs/} set
$\sigma = 0$.

\textbf{Configuration.} Benchmarks are built by
\texttt{benchmarks.build\_benchmark(config)}, which accepts both a flat layout
(\texttt{benchmark: <name>} plus top-level \texttt{noise\_std}/\texttt{bounds})
and a nested layout (\texttt{benchmark: \{name, bounds, noise\_std, \dots\}}).
Bounds and intrinsic parameters (e.g.\ \texttt{rate} for
\texttt{exponential\_1d}; \texttt{trap\_location}, \texttt{width},
\texttt{shape}, \texttt{trap\_gap} for \texttt{gaussian\_trap\_1d}) may be
overridden per benchmark; unknown keys are rejected so typos cannot be silently
ignored.

Each benchmark exposes the following entry points used by the pipeline:
\begin{description}
  \item[\texttt{f(X)}] : Noiseless evaluation. Used \emph{only} post-hoc for metrics; the BO algorithm never sees it.
  \item[\texttt{observe(X)}] : Noisy evaluation $f(X) + \mathcal{N}(0,\sigma^2)$. The only quantity the BO loop queries.
  \item[\texttt{test\_set(n, seed)}] : Deterministic scrambled Sobol grid in the native bounds, used for held-out (global) calibration. Reproducible from
        $(\text{benchmark}, n, \text{seed})$.
  \item[\texttt{near\_target\_test\_set(target, epsilon, n, seed)}] : Deterministic rejection-sampled Sobol grid of exactly $n$ points satisfying $|f(x)-T|\le\varepsilon$, used for the near-target calibration set (replaces masking the global grid).
  \item[\texttt{oracle\_target\_distance(T)}] : Analytic value of $\min_x|f(x)-T|$, currently implemented for \texttt{sine\_trend\_1d}; returns $0$ when the target is attainable. The default (other benchmarks) is \texttt{None}, treated downstream as $\texttt{oracle\_best}=0$.
\end{description}
 
% ============================================================
\subsection{Methods Compared}
\label{sec:methods}
% ============================================================
 
Two methods are registered in the pipeline under fixed symbolic names.
These names are the keys of every JSON file produced by all three pipeline stages.
 
\begin{center}
\begin{tabular}{lll}
\toprule
Symbolic key & Label & Role \\
\midrule
\texttt{right\_target\_matching\_ucb} & \RIGHT & SoC-compliant baseline \\
\texttt{wrong\_gp\_on\_abs\_lcb}      & \WRONG & SoC-violating baseline \\
\bottomrule
\end{tabular}
\end{center}
 
\subsubsection{WRONG: GP Trained on the Transformed Objective}
 
\textbf{Training targets:} $\tilde{y}_i = |y_i - T|$.
The GP surrogate family is held fixed across methods: like \RIGHT, \WRONG{}
uses a \texttt{FixedNoiseGP} at the known benchmark noise level $\sigma$, so the
\emph{only} difference between the two methods is the training target (raw $y$
vs.\ the transformed $|y-T|$), not the model class. (The noise on $|y-T|$ is in
truth non-Gaussian and folded near the target; fixing it to $\sigma$ keeps the
comparison controlled.)
 
\textbf{Acquisition:} negated lower confidence bound on the transformed quantity,
\[
  \alpha_t^{\mathrm{wrong}}(x)
  = -\mu_{g,t}(x) + \sqrt{\beta}\,\sigma_{g,t}(x),
\]
which BoTorch maximizes. The target $T$ is absorbed into the training data and
\emph{never appears explicitly in the acquisition class}.
 
\textbf{Key SoC violation:} the GP models $g = |f-T|$ rather than $f$, baking
task-specific structure into the surrogate.
 
\subsubsection{RIGHT: GP Trained on the Raw Function}
 
\textbf{Training targets:} $y_i = f(x_i) + \varepsilon_i$.
A \texttt{FixedNoiseGP} is used with the known benchmark noise level $\sigma$.
 
\textbf{Acquisition:} target-distance UCB on the raw GP,
\[
  \alpha_t^{\mathrm{right}}(x)
  = -|\mu_{f,t}(x) - T| + \sqrt{\beta}\,\sigma_{f,t}(x),
\]
which BoTorch maximizes.
The first term favors points whose posterior mean is close to $T$;
the second rewards uncertainty.
 
\textbf{SoC compliance:} the GP models $f$; the target $T$ enters only through the
acquisition function \texttt{TargetMatchingUCB}.
 
Both acquisitions are \texttt{AnalyticAcquisitionFunction} subclasses with $q=1$.
See \S\,\ref{sec:code-acq} for the implementation.
 
% ============================================================
\section{Experimental Protocol}
\label{sec:protocol}
% ============================================================
 
\subsection{Observation Noise}
 
All benchmarks use additive i.i.d.\ Gaussian noise.
The noise standard deviation $\sigma$ is fixed per benchmark and recorded in the YAML
configuration file (field \texttt{noise\_std}).
 
\subsection{Initial Design}
 
Each run starts from a Sobol initial design of size $n_0$.
The current \texttt{smoke\_1d.yaml} configuration sets $n_0 = 4$; the small-$n_0$
diagnostic configs use $n_0 = 1$.
 
\begin{contract}[Paired initial design]
\label{contract:paired}
For every seed, \RIGHT{} and \WRONG{} receive the \emph{same} Sobol initial design and
the \emph{same} initial noisy observations.
This is enforced by \texttt{run.py:\_run\_one\_seed}, which sets
\texttt{torch.manual\_seed(seed)} once, calls \texttt{problem.observe(X\_init)}, and
stores the result in \texttt{InitialDesign} at the seed level---\emph{outside} each
method's data structure.
\end{contract}
 
\subsection{Sequential Budget and Total History Length}
 
After initialization, each method runs for \texttt{budget} additional BO steps.
The full evaluation history of length $N = n_0 + \texttt{budget}$ is used for all
metrics. The current smoke configuration uses $\texttt{budget} = 50$ with
$n_0 = 4$, giving $N = 54$.
 
\subsection{Repetitions}
 
The current configurations use 100 independent seeds.
Each seed produces one independent draw of Sobol scramble, noise realizations, and
optimizer stochasticity.
 
\subsection{BO Iteration Loop}
 
At each BO iteration $t = 1, \ldots, \texttt{budget}$:
\begin{enumerate}
  \item Reconstruct the training set from all previous observations up to $t-1$,
        applying the method-specific transform (\texttt{"y\_obs"} or
        \texttt{"abs(y\_obs - target)"}).
  \item Fit the GP via marginal log-likelihood maximization (MLL).
  \item Record a \texttt{ModelSnapshot} (state dict of the fitted GP).
  \item Optimize the acquisition function to obtain the next candidate $x_t$.
  \item Query $y_t = \texttt{problem.observe}(x_t)$; record a \texttt{StepRecord}.
  \item Append $(x_t, y_t)$ to the history.
\end{enumerate}
 
\subsection{Configuration Reference}
 
\begin{lstlisting}[caption={Current smoke configuration (\texttt{smoke\_1d.yaml})},label=lst:config]
experiment:            exp_a_target_matching
tag:                   branin_2d_single_ring_noiseless_100seeds
benchmark:             branin_2d        # flat layout; or a nested mapping
bounds:                                 # optional per-run domain override
  lower:               [0.0, -2.0]
  upper:               [5.5,  6.0]
seeds:                 [1, 2, ..., 100]
target:                5.0
beta:                  2.0
initial_design_size:   4
budget:                50
noise_std:             0.0
num_restarts:          25
raw_samples:           2048
device:                cpu
methods:
  - right_target_matching_ucb
  - wrong_gp_on_abs_lcb

# Used by process.py and analyze.py only (not by run.py)
epsilon:               0.50       # hitting-time tolerance
failure_threshold:     1.00
test_grid:
  n_test:              4096       # global held-out grid size
  seed:                123
  calibration_every:   5          # calibration checkpoint cadence
  m_test_joint_distrib: 64        # joint-NLPD subgrid size (global)
  joint_nlpd_jitter:   1.0e-6
  # optional, for a dedicated near-target grid:
  # n_test_near_target: 512
  # near_target_seed:   124
  # m_test_joint_near_target: 64
near_target_epsilon:   0.50
calibration:
  compute_global:      true
  compute_near_target: true
  compute_joint:       true
  store_predictions:   true
aggregate:                        # consumed by analyze.py
  convergence_band_type:  pointwise   # pointwise | uniform | both
  calibration_band_type:  pointwise
  residual_band_type:     uniform
  n_bootstrap:            10000
  confidence:             0.95
  aggregate_calibration_curves:   true
  aggregate_instantaneous_curves: true
\end{lstlisting}
 
% ============================================================
\subsection{Quantities of Interest and Metrics}
\label{sec:quantities}
% ============================================================
 
All metrics operate on the \emph{noiseless} trajectory
$f_{1:N} = (f(x_1), \ldots, f(x_N))$, which covers the full evaluation history
including the initial design ($N = n_0 + \texttt{budget}$).
The noiseless values $f(x_t)$ are recorded during the BO loop and stored in
\texttt{StepRecord.f\_noiseless}.
The oracle value $\texttt{oracle\_best} = \min_x|f(x)-T|$ comes from the benchmark's
\texttt{oracle\_target\_distance(T)} method.
 
\subsubsection{Primary Per-Trajectory Curves}
\label{sec:curves}
 
The following four time series are stored in \texttt{ConvergenceCurves} (one per
(seed, method), length $N$):
 
\begin{center}
\begin{tabular}{lll}
\toprule
Field name & Symbol & Definition \\
\midrule
\texttt{instantaneous\_target\_distance}
  & $\dinst$
  & $|f(x_t) - T|$ \\[0.4em]
\texttt{running\_best\_target\_distance}
  & $\dbest$
  & $\min_{s \le t} |f(x_s) - T|$ \\[0.4em]
\texttt{cumulative\_inst\_target\_regret}
  & ---
  & $\sum_{s=1}^{t} \bigl(|f(x_s)-T| - \texttt{oracle\_best}\bigr)$ \\[0.4em]
\texttt{cumulative\_bsf\_target\_regret}
  & ---
  & $\sum_{s=1}^{t} \bigl(\dbest[s] - \texttt{oracle\_best}\bigr)$ \\
\bottomrule
\end{tabular}
\end{center}
 
\begin{remark}
  $\dinst$ is the \emph{instantaneous} target distance.
  It is \textbf{not} a simple regret: it does not accumulate knowledge across iterations.
  $\dbest$ is the running-best target distance;
  it is monotonically non-increasing.
\end{remark}
 
\subsubsection{Derived Scalars: Target Simple Regret}
 
The \textbf{target simple regret} at iteration $t$ is
\[
  \rtarget = \dbest - \texttt{oracle\_best}
           = \min_{s \le t}|f(x_s) - T| - \min_{x \in \mathcal{X}}|f(x) - T|.
\]
When the target is attainable ($\texttt{oracle\_best} = 0$), this simplifies to
$\rtarget = \dbest$.
For benchmarks where \texttt{oracle\_target\_distance} is not implemented,
\texttt{oracle\_best} is set to $0$ (conservative: attainability assumed).
 
\subsubsection{Per-Trajectory Scalars}
\label{sec:scalars}
 
The following scalars are stored in \texttt{FinalScalars}, computed by
\texttt{process.py} from the full history:
 
\begin{center}
\begin{longtable}{lp{8.5cm}}
\toprule
Field & Definition \\
\midrule
\endhead
\texttt{n\_evaluations}
  & $N = n_0 + \texttt{budget}$ \\
\texttt{final\_best\_target\_distance}
  & $\min_{1 \le t \le N} |f(x_t) - T|
    = d_N^{\mathrm{best}}$ \\
\texttt{final\_target\_simple\_regret}
  & $d_N^{\mathrm{best}} - \texttt{oracle\_best}$ \\
\texttt{mean\_inst\_target\_distance}
  & $\frac{1}{N}\sum_{t=1}^{N}|f(x_t)-T|$ \\
\texttt{median\_inst\_target\_distance}
  & $\mathrm{median}_{t}\,|f(x_t)-T|$ \\
\texttt{worst\_inst\_target\_distance}
  & $\max_{t}|f(x_t)-T|$ \\
\texttt{cumulative\_inst\_target\_regret}
  & $\sum_{t=1}^{N}(|f(x_t)-T|-\texttt{oracle\_best})$ \\
\texttt{cumulative\_bsf\_target\_regret}
  & $\sum_{t=1}^{N}(d_t^{\mathrm{best}}-\texttt{oracle\_best})$ \\
\texttt{auc\_running\_target\_regret}
  & Trapezoidal AUC of the running-best target regret
    curve $t \mapsto \rtarget$ \\
\texttt{hitting\_time}
  & $\min\{t : d_t^{\mathrm{best}} \le \varepsilon\}$;
    equals $N+1$ if never reached;
    \texttt{None} if $\varepsilon$ is unset \\
\texttt{failure\_indicator}
  & $\mathbf{1}\{d_N^{\mathrm{best}} > u\}$
    where $u =$ \texttt{failure\_threshold};
    \texttt{None} if threshold is unset \\
\texttt{best\_x}
  & $\arg\min_t |f(x_t)-T|$: the argmin point \\
\bottomrule
\end{longtable}
\end{center}
 
\subsubsection{Summary of All Metrics}
 
\begin{center}
\begin{tabular}{lll}
\toprule
Metric & Type & Primary use \\
\midrule
$\dinst = |f(x_t)-T|$ & curve & Visual exploration \\
$\dbest = \min_{s\le t}|f(x_s)-T|$ & curve & Convergence curve \\
$\rtarget = \dbest - \texttt{oracle\_best}$ & curve/scalar & Main perf.\ metric \\
Cumulative inst.\ regret & curve/scalar & Alternative perf.\ metric \\
Cumulative BSF regret & curve/scalar & Alternative perf.\ metric \\
AUC of running-best regret & scalar & Aggregate perf.\ \\
Hitting time $\tau_\varepsilon$ & scalar & Convergence speed \\
Failure indicator $\mathbf{1}\{\cdot > u\}$ & scalar & Tail-risk analysis \\
NLPD (marginal) & scalar (per checkpt) & Calibration quality \\
Joint Gaussian NLPD & scalar (per checkpt) & Joint calibration on a subgrid \\
Empirical 95\% coverage & scalar (per checkpt) & Calibration quality \\
Std.\ residual mean/std & scalar (per checkpt) & Calibration diagnostic \\
PIT mean/std & scalar (per checkpt) & Calibration diagnostic \\
Confidence violation rate & scalar (per checkpt) & Near-target calibration \\
\bottomrule
\end{tabular}
\end{center}
 
% ============================================================
\section{Experimental Output Pipeline}
\label{sec:pipeline}
% ============================================================
 
The pipeline is divided into three strictly separated stages:
 
\begin{stagebox}
\[
  \underbrace{\texttt{run.py}}_{\text{BO loop}}
  \;\longrightarrow\;
  \underbrace{\text{raw seed JSON}}_{\text{trajectory facts}}
  \;\longrightarrow\;
  \underbrace{\texttt{process.py}}_{\text{per-seed processing}}
  \;\longrightarrow\;
  \underbrace{\text{processed per-seed JSON}}_{\text{derived quantities}}
  \;\longrightarrow\;
  \underbrace{\texttt{analyze.py}}_{\text{cross-seed aggregation}}
  \;\longrightarrow\;
  \underbrace{\text{aggregate JSON}}_{\text{statistics}}
\]
\end{stagebox}
 
Each stage is a pure function of its inputs; only \texttt{run.py} touches the world
(draws noise, runs BO).
Re-running the BO loop is the only expensive operation; everything downstream is
recomputable in seconds.
 
\subsection{Run Directory Layout}
\label{sec:run-dir}
 
Each invocation of \texttt{run.py} creates one run directory:
 
\begin{lstlisting}
results/exp_a_target_matching/runs/<run_id>/
    config.json                   # YAML config saved as JSON
    metadata.json                 # git commit, sys.argv, Python, platform, UTC
    logs/run.log                  # streamed iteration log
    raw/<benchmark>/
        seed_001.json             # RAW: one file per seed
        seed_002.json
        ...
    processed/per_seed/<benchmark>/
        seed_001_processed.json   # PROCESSED: one file per seed
        seed_002_processed.json
        ...
    processed/aggregate/<benchmark>/
        final_performance.json    # AGGREGATE: up to 5 files per benchmark
        convergence_curves.json   #   performance / regret curves only
        calibration_summary.json  #   calibration scalars (final checkpoint)
        calibration_curves.json   #   calibration curves vs iteration (optional)
        statistical_tests.json
    figures/                      # Reserved (make_figures.py not yet implemented)
\end{lstlisting}
 
where \texttt{<run\_id>} has the form \texttt{YYYY-MM-DD\_HH-MM-SS\_<tag>} (UTC).
 
\subsubsection{RAW Data: \texttt{raw/<benchmark>/seed\_NNN.json}}
\label{sec:raw}
 
\paragraph{Produced by:} \texttt{run.py} (during the BO loop).
\paragraph{Granularity:} one file per (benchmark, seed).
\paragraph{Contents:} trajectory facts only 
 
\medskip
Top-level schema (\texttt{RawSeedData}, declared in \texttt{data/raw.py}):
\begin{lstlisting}
RawSeedData
|-- metadata         : SeedMetadata
|-- initial_design   : InitialDesign          # shared across methods
`-- methods          : {
        "right_target_matching_ucb": MethodTrajectory,
        "wrong_gp_on_abs_lcb"      : MethodTrajectory,
    }
\end{lstlisting}
 
\textbf{\texttt{SeedMetadata}:}
\texttt{experiment}, \texttt{problem\_type} (\texttt{"synthetic"}),
\texttt{benchmark}, \texttt{seed}, \texttt{budget}, \texttt{initial\_design\_size},
\texttt{target}, \texttt{noise\_std}, \texttt{dim}, \texttt{bounds}.
 
\textbf{\texttt{InitialDesign}:}
\texttt{X} (shape $(n_0, d)$), \texttt{y\_obs} (shape $(n_0,)$),
\texttt{f\_noiseless} (shape $(n_0,)$ or \texttt{None}).
Produced once per seed by \texttt{sobol\_design(bounds, n\_init, seed)}
followed by \texttt{problem.observe(X\_init)}.
 
\textbf{\texttt{MethodTrajectory}:}
method name, \texttt{train\_label} (transform recipe),
an ordered list of \texttt{budget} \texttt{StepRecord}s,
and a parallel list of \texttt{budget} \texttt{ModelSnapshot}s.
 
\textbf{\texttt{StepRecord} (per BO iteration):}
\begin{center}
\begin{tabular}{ll}
\toprule
Field & Content \\
\midrule
\texttt{iteration} & 1-based index (excludes initial design) \\
\texttt{x}         & candidate point, shape $(d,)$ \\
\texttt{y\_obs}    & noisy observation \\
\texttt{f\_noiseless} & noiseless benchmark value (for post-hoc metrics) \\
\texttt{acq\_value} & acquisition value at $x_t$ \\
\texttt{train\_size} & $n_0 + \texttt{iteration} - 1$ \\
\texttt{posterior\_mean\_at\_candidate} & $\mu_{t-1}(x_t)$ \\
\texttt{posterior\_std\_at\_candidate}  & $\sigma_{t-1}(x_t)$ \\
\bottomrule
\end{tabular}
\end{center}
 
\textbf{\texttt{ModelSnapshot} (per fitted GP):}
\texttt{iteration}, \texttt{model\_type} (\texttt{"SingleTaskGP"} or
\texttt{"FixedNoiseGP"}), \texttt{kernel\_name},
\texttt{noise\_std} (for \texttt{FixedNoiseGP} reconstruction;
\texttt{None} for inferred-noise models),
\texttt{state\_dict} (JSON-serialised \texttt{torch.nn.Module.state\_dict()},
including raw lengthscales, outputscale, input-normalisation min/range,
outcome-standardisation mean/std),
\texttt{hyperparameters} (parsed scalars for human inspection only).
 
\textbf{What the raw file does NOT store:}
$(\texttt{train\_X}, \texttt{train\_Y})$ at each iteration (rebuilt by
\texttt{process.py}), aggregated quantities, and test-grid predictions.
 
\subsubsection{PROCESSED Per-Seed Data:
  \texttt{processed/per\_seed/<benchmark>/seed\_NNN\_processed.json}}
\label{sec:processed}
 
\paragraph{Produced by:} \texttt{process.py} (post-processing; reads one raw file at a time).
\paragraph{Granularity:} one file per (benchmark, seed).
\paragraph{Contents:} all derived quantities for that seed --- scalars, convergence
curves, and calibration reports for every retained checkpoint.
 
\medskip
Top-level schema (\texttt{ProcessedSeedData}, declared in \texttt{data/processed.py}):
\begin{lstlisting}
ProcessedSeedData
|-- metadata          : dict           # copy of SeedMetadata (no raw lookup needed)
|-- target            : float
|-- epsilon           : float | None
|-- failure_threshold : float | None
|-- oracle_best       : float          # min_x |f(x) - T|
|-- test_grid         : TestGrid        # see below; stored once per seed
`-- methods           : { "<name>": MethodProcessed }

TestGrid                                # grid provenance, no points stored
|-- n_test, seed                        # global Sobol grid
|-- near_target_epsilon, near_target_mode
|-- n_test_near_target, near_target_seed, near_target_pool_size
`-- m_test_joint_distrib, m_test_joint_near_target

MethodProcessed
|-- method_name   : str
|-- train_label   : str                # "y_obs" or "abs(y_obs - target)"
|-- final         : FinalScalars       # see Section 5.3
|-- convergence   : ConvergenceCurves  # see Section 5.1
`-- calibrations  : list[CalibrationReport]
\end{lstlisting}

Test-grid metadata (sizes, seeds, near-target/joint settings) is stored once
per seed in \texttt{ProcessedSeedData.test\_grid} and is not duplicated inside
each report. Each \texttt{CalibrationReport} contains:
\begin{itemize}
  \item \texttt{iteration\_index}: the checkpoint iteration index;
  \item \texttt{global\_calibration}: full-grid diagnostic dictionary
        (\texttt{calibration\_summary} on $X_{\mathrm{test}}$);
  \item \texttt{near\_target} (optional): diagnostics on the dedicated
        near-target grid, with \texttt{near\_target\_count}/\texttt{fraction}
        and, when configured, joint near-target NLPD fields;
  \item \texttt{joint\_posterior} (optional): joint Gaussian NLPD on an
        $m_{\mathrm{joint}}$-sized subgrid (Mahalanobis term, log-det, jitter);
  \item \texttt{predictions} (optional): \texttt{CalibrationPredictionGrid}
        storing only $(\mu,\sigma)$ on the global grid --- $X_{\mathrm{test}}$ and
        $y_{\mathrm{true}}$ are reconstructed deterministically from the
        benchmark, so they are no longer duplicated per checkpoint.
\end{itemize}
 
\subsubsection{AGGREGATED Cross-Seed Data:
  \texttt{processed/aggregate/<benchmark>/*.json}}
\label{sec:aggregate}
 
\paragraph{Produced by:} \texttt{analyze.py} (reads \emph{all} processed per-seed files for a benchmark; the only stage that opens more than one seed file).
\paragraph{Granularity:} one set of up to five files per benchmark
(\texttt{calibration\_curves.json} is written only when
\texttt{aggregate.aggregate\_calibration\_curves} is true).

 
\subsubsection{\texttt{final\_performance.json}}
 
For each method and each scalar in
\[
  \texttt{\_FINAL\_SCALAR\_KEYS} = \{
    \texttt{final\_best\_target\_distance},\;
    \texttt{final\_target\_simple\_regret},\;
    \texttt{mean\_inst\_target\_distance},
\]
\[
    \texttt{worst\_inst\_target\_distance},\;
    \texttt{cumulative\_inst\_target\_regret},\;
    \texttt{cumulative\_bsf\_target\_regret},\;
    \texttt{auc\_running\_target\_regret},\;
    \texttt{hitting\_time}
  \},
\]
the per-seed values are stacked (dropping \texttt{None}) and summarised by
\texttt{common.statistics.summarize\_sample}, which returns:
\[
  \{n,\;\mathrm{mean},\;\mathrm{std},\;\mathrm{stderr},\;\mathrm{median},\;
    q_{25},\;q_{75},\;\min,\;\max,\;q_{90},\;q_{95}\}.
\]
When \texttt{failure\_indicator} is available, the empirical failure rate is also
reported.
 
\subsubsection{\texttt{convergence\_curves.json}}

This file holds the \emph{performance/regret} curves only (calibration curves
live in \texttt{calibration\_curves.json}). For each method and each of the four
time-series (\S\,\ref{sec:curves}), the per-seed curves are stacked into a matrix
of shape $(n_{\mathrm{seeds}}, N)$ and aggregated by
\texttt{aggregate\_convergence\_curves}, producing
\[
  \{\texttt{center},\;\texttt{lower},\;\texttt{upper},\;\texttt{band\_type}\}
\]
per (method, curve type). The band is selected by
\texttt{aggregate.convergence\_band\_type} (default \texttt{pointwise}):
\begin{itemize}
  \item \texttt{pointwise}: per-step percentile bootstrap CI
        (\texttt{bootstrap\_ci\_curves}, \S\,\ref{sec:bootstrap});
  \item \texttt{uniform}: simultaneous (Kolmogorov-style) band around the chosen
        statistic (\texttt{bootstrap\_uniform\_band}), also reporting the
        half-width multiplier $q$;
  \item \texttt{both}: emits a \texttt{pointwise} and a \texttt{uniform} sub-block.
\end{itemize}
The central statistic (median by default, mean optional) is shared by both band
types.

\subsubsection{\texttt{calibration\_summary.json}}

For each method, \texttt{summarize\_sample} is applied to the per-seed values of each
calibration scalar (NLPD, coverage, residual statistics, PIT, violation rate,
and---when computed---joint Gaussian NLPD), restricted to the \emph{final}
calibration checkpoint of each seed. Three buckets are reported when available:
\texttt{global} (full grid), \texttt{near\_target} (dedicated near-target grid,
plus \texttt{near\_target\_count}/\texttt{fraction} and joint near-target NLPD),
and \texttt{joint} (global joint-posterior NLPD).

\subsubsection{\texttt{calibration\_curves.json} (optional)}

Written only when \texttt{aggregate.aggregate\_calibration\_curves} is true.
It is the calibration analogue of \texttt{convergence\_curves.json}: for each
method and each calibration scalar, the per-checkpoint values are stacked across
seeds and aggregated into a band \emph{versus iteration}. Band type is chosen
\emph{per metric}: standardized-residual metrics use
\texttt{aggregate.residual\_band\_type} (default \texttt{uniform}); all other
calibration curves (NLPD, coverage, violation rate, PIT, joint posterior) use
\texttt{aggregate.calibration\_band\_type} (default \texttt{pointwise}).

\subsubsection{\texttt{statistical\_tests.json}}

For every (\RIGHT, \WRONG) pair matched by the \texttt{right\_*}/\texttt{wrong\_*}
prefixes, and for every scalar in \texttt{\_FINAL\_SCALAR\_KEYS}, the file reports
the \textbf{paired Wilcoxon signed-rank test as the primary comparison} (the
seeds are paired by construction; \S\,\ref{sec:protocol}), alongside the
unpaired Mann--Whitney $U$ test and the effect sizes:
\[
  \bigl\{
    n_R,\;n_W,\;
    W_{\mathrm{Wilcoxon}},\;p_{\mathrm{Wilcoxon}},\;
    U_{\mathrm{MW}},\;p_{\mathrm{MW}},\;
    \rho_{\mathrm{rb}},\;
    d_{\mathrm{Cohen}},\;
    \text{failure rates}
  \bigr\}.
\]
See \S\,\ref{sec:stats} for definitions.
 
\medskip
\textbf{Summary of the pipeline:}
 
\begin{center}
\begin{tabular}{lllll}
\toprule
Stage & Script & Per-seed or agg. & During BO? & Downstream use \\
\midrule
Raw & \texttt{run.py} & per seed & yes & \texttt{process.py} \\
Processed & \texttt{process.py} & per seed & no & \texttt{analyze.py}, notebook \\
Aggregate perf.\ & \texttt{analyze.py} & aggregate & no & notebook, figures \\
Aggregate curves & \texttt{analyze.py} & aggregate & no & notebook, figures \\
Aggregate calib.\ & \texttt{analyze.py} & aggregate & no & notebook, figures \\
Aggregate calib.\ curves & \texttt{analyze.py} & aggregate (opt.) & no & notebook, figures \\
Statistical tests & \texttt{analyze.py} & aggregate & no & notebook, paper \\
\bottomrule
\end{tabular}
\end{center}
 
% ============================================================
\subsection{Calibration Pipeline}
\label{sec:calibration}
% ============================================================
 
\subsubsection{Purpose}
 
The calibration pipeline answers: does the GP's uncertainty quantification remain valid
across the search space, and in particular near the target level set?
It is the empirical bridge between model quality and optimization behavior.
 
\subsubsection{Held-Out Test Grid}
 
A deterministic global Sobol grid $X_{\mathrm{test}}$ (size $n_{\mathrm{test}}$,
e.g.\ $4096$ in \texttt{smoke\_1d.yaml}) is produced by
\texttt{benchmark.test\_set(n\_test, seed)}. When
\texttt{near\_target\_epsilon} and \texttt{n\_test\_near\_target} are set, a
\emph{dedicated} near-target grid $X_{\mathrm{near}}$ of exactly
$n_{\mathrm{near}}$ points satisfying $|f(x)-T|\le\varepsilon$ is produced by
\texttt{benchmark.near\_target\_test\_set(target, epsilon, n, seed)}
(rejection-sampled Sobol). This replaces the previous approach of masking the
global grid, guaranteeing a controlled near-target sample size even when the
band is thin.

\begin{contract}[Shared deterministic test grids]
\label{contract:testgrid}
$X_{\mathrm{test}}$ and $X_{\mathrm{near}}$ are pure functions of the benchmark
and their respective $(\text{size},\text{seed})$.
The same grids are used for both methods and across all BO seeds within a run.
\end{contract}
 
\subsubsection{GP Reconstruction at Each Checkpoint}
 
Calibration is evaluated at every iteration that is a multiple of
\texttt{calibration\_every} (default: 5) and at the final iteration.
At each retained \texttt{ModelSnapshot}:
 
\begin{enumerate}
  \item \textbf{Rebuild training data:} reconstruct
        $(\texttt{train\_X}, \texttt{train\_Y})$ from the initial design and the
        trajectory prefix up to that iteration, applying the method's
        \texttt{train\_label} transform.
  \item \textbf{Reconstruct GP (no MLL refitting):}
        \begin{enumerate}[label=(\alph*)]
          \item Call \texttt{GPModelBuilder.construct(X, Y)} to create the GP object.
          \item Call \texttt{load\_state\_dict\_into(model, snapshot.state\_dict)}.
        \end{enumerate}
        The result is bit-identical to the GP used during BO (Contract~\ref{contract:no-refit}).
  \item \textbf{Posterior on test grid:} evaluate $(\mu, \sigma) =
        \texttt{GPModelBuilder.predict}(\texttt{model}, X_{\mathrm{test}})$.
  \item \textbf{Ground truth:}
        $y_{\mathrm{true}} = f(X_{\mathrm{test}})$ for \RIGHT;
        $y_{\mathrm{true}} = |f(X_{\mathrm{test}}) - T|$ for \WRONG{} (method-aware).
  \item \textbf{Full-grid calibration:} compute the summary in \S\,\ref{sec:calib-diag}
        on $X_{\mathrm{test}}$ (key \texttt{global\_calibration}).
  \item \textbf{Near-target calibration} (if configured): compute the same summary
        on the dedicated near-target grid $X_{\mathrm{near}}$ (key
        \texttt{near\_target}), recording \texttt{near\_target\_count} and
        \texttt{fraction}. The near-target band is always defined against raw $f$,
        regardless of what the GP models.
  \item \textbf{Joint Gaussian NLPD} (if \texttt{compute\_joint} and
        \texttt{m\_test\_joint\_distrib} are set): materialise the full posterior
        covariance over the first $m_{\mathrm{joint}}$ test points and compute the
        Cholesky-stabilised joint NLPD (key \texttt{joint\_posterior}); optionally
        also on the near-target slice.
  \item \textbf{Store prediction grid} (if \texttt{store\_predictions}): save
        only $(\mu, \sigma)$ on $X_{\mathrm{test}}$ as
        \texttt{CalibrationPredictionGrid}. $X_{\mathrm{test}}$ and
        $y_{\mathrm{true}}$ are reconstructed deterministically downstream, so
        they are no longer stored per checkpoint.
\end{enumerate}
 
\begin{contract}[Calibration uses load-only reconstruction]
\label{contract:no-refit}
\texttt{process.py} rebuilds GPs exclusively via
\texttt{construct} + \texttt{load\_state\_dict\_into}.
MLL is \textbf{never} refitted during processing.
\texttt{load\_state\_dict\_into} is strict: it raises on key mismatch.
\end{contract}
 
\subsubsection{Calibration Diagnostics}
\label{sec:calib-diag}
 
Let $y_j$ be the ground truth, $\mu_j$ the posterior mean, and $\sigma_j$ the posterior
standard deviation at test point $j$.
Define the \textbf{standardized residual}
\[
  Z_j = \frac{y_j - \mu_j}{\sigma_j}.
\]
Under correct calibration, $Z_j \approx \mathcal{N}(0,1)$.
The \textbf{probability integral transform (PIT)} value is
\[
  U_j = \Phi(Z_j),
\]
which should be approximately $\mathrm{Uniform}(0,1)$ under correct calibration.
 
\begin{center}
\begin{tabular}{llp{6.5cm}}
\toprule
Key & Symbol / formula & Description \\
\midrule
\texttt{nlpd}
  & $-\frac{1}{m}\sum_j \log p(y_j \mid \mathcal{D}_t)$
  & Negative log predictive density \\
\texttt{coverage\_95}
  & $\frac{1}{m}\sum_j \mathbf{1}\{|Z_j|\le 1.96\}$
  & Empirical 95\% interval coverage \\
\texttt{residual\_mean}
  & $\bar{Z}$
  & Mean standardized residual (should be $\approx 0$) \\
\texttt{residual\_std}
  & $\mathrm{std}(Z_j)$
  & Std of residuals (should be $\approx 1$) \\
\texttt{residual\_mean\_abs}
  & $\frac{1}{m}\sum_j |Z_j|$
  & Mean absolute residual \\
\texttt{pit\_mean}
  & $\bar{U}$
  & Mean PIT (should be $\approx 0.5$) \\
\texttt{pit\_std}
  & $\mathrm{std}(U_j)$
  & Std of PIT (should be $\approx 1/\sqrt{12}$) \\
\texttt{confidence\_violation\_rate}
  & $\frac{1}{m}\sum_j \mathbf{1}\{|f(x_j)-\mu_j| > \sqrt{\beta}\,\sigma_j\}$
  & Fraction violating UCB confidence bound at $\beta = 2$ \\
\texttt{nlpd\_joint}
  & $\frac{1}{2m}\bigl[(y{-}\mu)^\top K^{-1}(y{-}\mu) + \log\det K + m\log 2\pi\bigr]$
  & Joint Gaussian NLPD over an $m$-point subgrid (full posterior covariance $K$;
    Cholesky with jitter). Reported with \texttt{mahalanobis\_sq} and \texttt{logdet}. \\
\bottomrule
\end{tabular}
\end{center}
 
The near-target bucket additionally reports
\texttt{near\_target\_count} and \texttt{near\_target\_fraction}.
 
\subsubsection{Theoretical Motivation for Near-Target Calibration}
 
UCB-type BO convergence relies on the uniform confidence event
\[
  \mathcal{E}
  = \bigl\{
    \forall x,\,\forall t,\;
    |f(x) - \mu_{t-1}(x)| \le \sqrt{\beta_t}\,\sigma_{t-1}(x)
  \bigr\}.
\]
A SoC violation can invalidate $\mathcal{E}$ specifically near the target.
If the GP models $g = |f-T|$, its posterior variance $\sigma_{g,t}^2(x)$ may be
small near kinks---exactly where $f$ crosses $T$---even though uncertainty about
$f$ near the target is high.
The near-target confidence violation rate
\[
  V_t^{\mathrm{near}}
  = \frac{1}{|\mathcal{X}_{\mathrm{near}}|}
    \sum_{x_j \in \mathcal{X}_{\mathrm{near}}}
    \mathbf{1}\bigl\{|f(x_j) - \mu_t(x_j)| > \sqrt{\beta}\,\sigma_t(x_j)\bigr\}
\]
directly probes this mechanism.
 
% ============================================================
\section{Statistical Analysis}
\label{sec:stats}
% ============================================================
 
\subsection{Descriptive Summaries}
 
For each (method, scalar) pair, \texttt{final\_performance.json} contains
$\{n, \mathrm{mean}, \mathrm{std}, \mathrm{stderr}, \mathrm{median}, q_{25}, q_{75}, \min, \max, q_{90}, q_{95}\}$ computed across seeds. These descriptive summaries are the first thing to inspect in the notebook.
 
\subsection{Bootstrap Confidence Intervals for Convergence Curves}
\label{sec:bootstrap}
 
\textbf{What is resampled:} the $n_{\mathrm{seeds}}$ rows of the curve matrix
(shape $n_{\mathrm{seeds}} \times N$).
Each row is one independent seed; resampling rows captures seed-to-seed variability.
 
\textbf{What is estimated:} the sampling distribution of the chosen central tendency
(median by default, mean optional) of the curve at each fixed iteration $t$.
 
\textbf{Procedure:} draw $B = 5000$ bootstrap replicates; at each iteration $t$,
record the statistic across resampled seeds; report $q_{2.5\%}$ and $q_{97.5\%}$
as the lower and upper bounds. No parametric model is assumed for the per-iteration metric, which matters because
target-distance distributions are bounded below by zero and may be heavily skewed.
 
\begin{remark}
The resulting CI is \textbf{pointwise} in $t$.
The probability that the entire true curve lies simultaneously inside the envelope
is strictly smaller than $1 - \alpha$.
These bands are visual uncertainty summaries, not uniform-in-time confidence bands.
\end{remark}
 
\subsection{Paired Wilcoxon Signed-Rank Test (primary)}

For each scalar in \texttt{\_FINAL\_SCALAR\_KEYS}, the \emph{primary} test is a
paired Wilcoxon signed-rank test
(\texttt{common.statistics.wilcoxon\_paired\_comparison}) on the per-seed
differences $\Delta_i = w_i - r_i$, which exploits the paired-seed design
(\S\,\ref{sec:protocol}):
\begin{itemize}
  \item \texttt{alternative = "greater"}: a small $p$-value supports
        $\mathrm{median}(\Delta) > 0$, i.e.\ \WRONG{} has larger regret than
        \RIGHT{} (\RIGHT{} is better, smaller-is-better convention).
  \item Non-parametric; uses the within-seed pairing rather than treating the
        two samples as independent.
  \item Reported with the median/mean of $\Delta$ (\texttt{wrong} minus
        \texttt{right}). Degenerate all-zero-difference cases return \texttt{NaN}
        with an explanatory flag rather than raising.
\end{itemize}

\subsection{Mann--Whitney Test (secondary)}

For each scalar in \texttt{\_FINAL\_SCALAR\_KEYS}, a Mann--Whitney U test
is also run between the \RIGHT{} and \WRONG{} samples as an unpaired cross-check:
\begin{itemize}
  \item \texttt{alternative = "less"}: a small $p$-value supports the one-sided
        hypothesis $r_N^{\mathrm{right}} < r_N^{\mathrm{wrong}}$ (RIGHT is better).
  \item The test is non-parametric and does not assume normality.
  \item It is reported alongside effect sizes; $p$-value alone is insufficient.
\end{itemize}
 
\subsection{Effect Sizes}
 
\paragraph{Rank-biserial correlation.}
\[
  \rho_{\mathrm{rb}}
  = \frac{\#\{(r,w): w > r\} - \#\{(r,w): w < r\}}{n_R\,n_W}
  \in [-1,1].
\]
Positive values mean \RIGHT{} tends to have smaller values than \WRONG{}.
Ties are treated neutrally. Computed exactly via vectorised pairwise comparison.
 
\paragraph{Cohen's $d$ (right-better convention).}
\[
  d = \frac{\bar{w} - \bar{r}}{s_{\mathrm{pooled}}},
  \qquad
  s_{\mathrm{pooled}}^2
  = \frac{(n_R-1)s_r^2 + (n_W-1)s_w^2}{n_R + n_W - 2}.
\]
Positive values mean \RIGHT{} is better.
Returns \texttt{NaN} when either sample has fewer than two points or the pooled
standard deviation is effectively zero. Should be interpreted with caution for
skewed or outlier-heavy distributions.
 
\subsection{Failure-Rate Comparison}
 
When \texttt{failure\_threshold} $u$ is set, the empirical failure probability
$\hat{p}_u = \frac{1}{n}\sum_i \mathbf{1}\{d_{N,i}^{\mathrm{best}} > u\}$
is reported for each method and included in \texttt{statistical\_tests.json}.
This targets the tail of the performance distribution rather than its center.
 
\subsection{Pairing}
 
The seeds are technically \emph{paired}: \RIGHT{} and \WRONG{} share the same Sobol
initial design and the same \texttt{torch.manual\_seed} call
(Contract~\ref{contract:paired}).
The paired Wilcoxon signed-rank test (now the primary comparison) exploits this
structure; the Mann--Whitney test is reported as an unpaired secondary check.

\begin{remark}
\texttt{analyze.py} builds seed-for-seed matched samples
(\texttt{\_stack\_scalars\_paired}) so the Wilcoxon test only consumes index-for-index
paired values. The previous ``planned'' status of this test is now resolved:
\texttt{wilcoxon\_paired\_comparison} is implemented in
\texttt{common/statistics.py}.
\end{remark}
 
\subsection{Sign Convention Summary}
 
\begin{center}
\begin{tabular}{lll}
\toprule
Quantity & Positive means & Implementation \\
\midrule
$\rho_{\mathrm{rb}}$ & \RIGHT{} beats \WRONG{} & exact pairwise \\
Cohen's $d$ & \RIGHT{} beats \WRONG{} & $\bar{w} - \bar{r}$ / pooled std \\
MW \texttt{alternative} & \RIGHT{} $<$ \WRONG{} & \texttt{"less"} \\
\bottomrule
\end{tabular}
\end{center}
 
% ============================================================
\section{Notebook Analysis Roadmap}
\label{sec:notebook}
% ============================================================
 
This section provides a step-by-step guide for writing the analysis notebook.
Each step specifies which output files to read.
 
\subsection{Step 0: Setup and Paths}
 
\begin{itemize}
  \item Set \texttt{run\_dir} to the target run directory
        \texttt{results/exp\_a\_target\_matching/runs/<run\_id>/}.
  \item Set \texttt{benchmark} to the benchmark name (e.g., \texttt{branin\_2d}).
  \item All aggregate files are under
        \texttt{<run\_dir>/processed/aggregate/<benchmark>/}.
  \item All per-seed files are under
        \texttt{<run\_dir>/processed/per\_seed/<benchmark>/}.
\end{itemize}
 
\subsection{Step 1: Load and Inspect Configuration}
 
\begin{itemize}
  \item \textbf{Read:} \texttt{<run\_dir>/config.json},
        \texttt{<run\_dir>/metadata.json}.
  \item Verify: benchmark, target, $n_0$, budget, seeds, methods, noise level.
  \item Check provenance: git commit, UTC timestamp.
\end{itemize}
 
\subsection{Step 2: Final Performance Summary}
 
\begin{itemize}
  \item \textbf{Read:} \texttt{processed/aggregate/<benchmark>/final\_performance.json}.
  \item Display descriptive summary table (mean, median, IQR) for
        \texttt{final\_best\_target\_distance} and
        \texttt{final\_target\_simple\_regret} for each method.
  \item Check failure rates and hitting-time distributions if available.
  \item Cross-reference with
        \texttt{processed/aggregate/<benchmark>/statistical\_tests.json}
        for $p$-values and effect sizes ($\rho_{\mathrm{rb}}$, Cohen's $d$).
\end{itemize}
 
\subsection{Step 3: Convergence Curves}
 
\begin{itemize}
  \item \textbf{Read:} \texttt{processed/aggregate/<benchmark>/convergence\_curves.json}.
  \item Plot the median running-best target regret $\rtarget$ against iteration $t$
        for both methods, with bootstrap CI bands.
  \item Optionally overlay mean curves (from the same file if \texttt{mean} was computed).
  \item Note: the curve matrix per (method, curve type) is shape
        $(n_{\mathrm{seeds}}, N)$; individual seed curves may be overlaid for
        qualitative inspection.
\end{itemize}
 
\subsection{Step 4: Calibration Analysis}
 
\begin{itemize}
  \item \textbf{Read:} \texttt{processed/aggregate/<benchmark>/calibration\_summary.json}
        for cross-seed aggregate view.
  \item Plot NLPD trajectory (per checkpoint) for both methods.
  \item Compare empirical 95\% coverage and confidence violation rate.
  \item For per-seed detail, read individual
        \texttt{processed/per\_seed/<benchmark>/seed\_NNN\_processed.json}:
        access \texttt{methods["<name>"]["calibrations"][-1]} for the final checkpoint.
  \item From \texttt{CalibrationPredictionGrid}: plot PIT histogram, QQ-plot of
        standardized residuals.
  \item Compare full-grid vs.\ near-target diagnostics; report
        \texttt{near\_target\_fraction} to understand how much of the test grid is
        near the target.
\end{itemize}
 
\subsection{Step 5: Failure and Tail-Risk Analysis}
 
\begin{itemize}
  \item \textbf{Read:} \texttt{final\_performance.json} (high quantiles
        $q_{90}$, $q_{95}$, worst run).
  \item Identify seeds where one method failed (\texttt{failure\_indicator = 1}).
  \item For representative failing seeds, load the per-seed processed file to inspect
        convergence curves and calibration trajectories.
  \item Compare $\hat{p}_u$ between \RIGHT{} and \WRONG{}.
\end{itemize}
 
\subsection{Step 6: Representative Seed Inspection (1D benchmark)}
 
\begin{itemize}
  \item \textbf{Read:} one or more
        \texttt{processed/per\_seed/sine\_trend\_1d/seed\_NNN\_processed.json}.
  \item From \texttt{CalibrationPredictionGrid}: reconstruct posterior mean/variance
        plots at selected iterations (e.g., $t \in \{5, 10, 20, 40\}$) for both methods.
  \item Plot $f$, $|f-T|$, posterior mean, and 95\% credible intervals side by side
        (Figure A1/A2 basis).
\end{itemize}
 
\subsection{Step 7: Export Figures}
 
\begin{itemize}
  \item Save all figures to \texttt{<run\_dir>/figures/} with descriptive names.
  \item Record which file, method, and seeds each figure uses.
  \item \textit{Note:} a \texttt{make\_figures.py} script for publication-quality
        figure generation is \textbf{not yet implemented}; the notebook is the current
        primary analysis tool.
\end{itemize}
 
\subsection{Notebook Dependency Map}
 
\begin{center}
\begin{tabular}{lll}
\toprule
Notebook step & Primary file(s) to read & Secondary file(s) \\
\midrule
Config / provenance & \texttt{config.json}, \texttt{metadata.json} & --- \\
Final performance & \texttt{final\_performance.json} & \texttt{statistical\_tests.json} \\
Convergence curves & \texttt{convergence\_curves.json} & per-seed files (optional) \\
Calibration summary & \texttt{calibration\_summary.json} & per-seed files \\
Failure / tail risk & \texttt{final\_performance.json} & per-seed files \\
1D visual inspection & per-seed files & raw files (optional) \\
\bottomrule
\end{tabular}
\end{center}
 
% ============================================================
\section{Figures and Tables to Produce}
\label{sec:figures}
% ============================================================
 
The \texttt{figures/} subdirectory is reserved for publication-quality outputs.
A \texttt{make\_figures.py} script is \textbf{not yet implemented}; figures are
currently generated in the analysis notebook.
 
\subsection{Figure A1 --- 1D Mechanism Figure}
 
On \texttt{sine\_trend\_1d}: plot $f$, the target level $T$, the transformed function
$|f-T|$, the \WRONG{} GP posterior on $|f-T|$, the \RIGHT{} GP posterior on $f$, and
sampled points at a representative iteration.
\textbf{Goal:} make visually clear that $|f-T|$ contains artificial kinks even when $f$
is smooth.
\textbf{Source:} \texttt{CalibrationPredictionGrid} from a per-seed processed file.
 
\subsection{Figure A2 --- Posterior Snapshots}
 
For selected iterations $t \in \{5, 10, 20, 40\}$ on \texttt{sine\_trend\_1d}:
plot posterior means and 95\% credible intervals for both methods, reconstructed offline
from \texttt{ModelSnapshot} state dicts via \texttt{load\_state\_dict\_into}.
\textbf{Source:} \texttt{CalibrationPredictionGrid} from a per-seed processed file.
 
\subsection{Figure A3 --- Convergence Curves}
 
Median running-best target regret $t \mapsto r_t^{\mathrm{target}}$ over all seeds,
with 95\% pointwise bootstrap CI.
Both methods overlaid on the same axes.
\textbf{Source:} \texttt{convergence\_curves.json}.
 
\subsection{Figure A4 --- Calibration Diagnostics}
 
For each method:
\begin{itemize}
  \item NLPD trajectory across calibration checkpoints;
  \item empirical coverage trajectory;
  \item PIT histogram at the final checkpoint;
  \item QQ-plot of standardized residuals at the final checkpoint;
  \item near-target confidence violation rate.
\end{itemize}
\textbf{Source:} \texttt{calibration\_summary.json} (aggregate view);
per-seed processed files for individual-seed PIT/QQ plots.
 
\subsection{Summary Performance Table}
 
A table with rows = methods, columns = descriptive statistics
(median $\pm$ IQR, mean $\pm$ SE, $q_{90}$, $q_{95}$, failure rate) for
\texttt{final\_best\_target\_distance} and \texttt{final\_target\_simple\_regret}.
Effect sizes ($\rho_{\mathrm{rb}}$, Cohen's $d$) and $p$-values in a companion table.
\textbf{Source:} \texttt{final\_performance.json}, \texttt{statistical\_tests.json}.
 
% ============================================================
\section{Code Architecture}
\label{sec:code}
% ============================================================
 
\subsection{Design Principles}
 
Three principles organize the codebase:
\begin{enumerate}
  \item \textbf{Problem abstraction:} the BO loop talks to a generic
        \texttt{ProblemSpec}; it does not know whether observations come from an
        analytic benchmark, a tabular dataset, or a real experiment.
  \item \textbf{Shared infrastructure:} GP construction, metrics, diagnostics,
        statistics, logging, and serialization are implemented once in \texttt{common/}
        and reused across experiments.
  \item \textbf{Task-specific acquisitions:} acquisition functions encoding a particular
        SoC scenario live in \texttt{acquisitions/}.
\end{enumerate}
 
The dependency direction is strictly unidirectional:
\[
  \texttt{data/}
  \;\subset\;
  \texttt{common/}
  \;\subset\;
  \texttt{benchmarks/},\;\texttt{acquisitions/}
  \;\subset\;
  \texttt{experiments/}.
\]
 
\subsection{Directory Layout}
 
\begin{lstlisting}
bayesian-optimization-separation-concerns/
|-- common/
|   |-- problems.py          # ProblemSpec abstract base + Observation
|   |-- models.py            # GP construction, fitting, state-dict I/O
|   |-- initial_design.py    # Sobol initial design
|   |-- acquisition_utils.py # Wrappers around optimize_acqf
|   |-- metrics.py           # target-distance / regret / hitting-time
|   |-- diagnostics.py       # NLPD, coverage, PIT, residuals, near-target
|   |-- statistics.py        # Across-seed aggregation + pairwise tests
|   |-- logging.py           # RunLogger: run directory lifecycle
|   `-- serialization.py     # Tensor/JSON/CSV/Pickle helpers
|
|-- benchmarks/
|   |-- synthetic.py         # SyntheticBenchmark base + near_target_test_set
|   |-- toy.py               # 8 benchmarks (4 standard + 4 diagnostic 1D)
|   `-- __init__.py          # BENCHMARK_REGISTRY + build_benchmark(config)
|
|-- acquisitions/
|   `-- target_matching.py   # TargetMatchingUCB + WrongTargetMatchingUCB
|
|-- data/
|   |-- raw.py               # RawSeedData schema
|   `-- processed.py         # ProcessedSeedData, TestGrid, CalibrationReport
|
|-- experiments/
|   `-- exp_a_target_matching/
|       |-- run.py           # Stage 1: generates raw seed JSON
|       |-- process.py       # Stage 2: builds processed per-seed JSON
|       |-- analyze.py       # Stage 3: builds aggregate JSON
|       |-- configs/         # 12 YAML configs (one per benchmark / variant)
|       `-- notebooks/       # Exp_A_analysis_{sine,forrester,branin,hartmann}.ipynb
|
|-- scripts/                 # run_all_cpu.sh, process_last_runs.sh, smoke wrapper
|-- hpc/narval/              # Reserved for SLURM job scripts (currently empty)
|-- results/                 # One subtree per experiment / run
|-- CLAUDE.md, README.md, requirements.txt
\end{lstlisting}
 
\subsection{\texttt{common/} --- Reusable Infrastructure}
 
\paragraph{\texttt{common/problems.py}.}
\begin{itemize}
  \item \texttt{Observation}: dataclass $(X,\;y_{\mathrm{obs}},\;f_{\mathrm{noiseless}},\;
        \mathrm{metadata})$. Shapes validated in \texttt{\_\_post\_init\_\_}.
  \item \texttt{ProblemSpec(ABC)}: holds \texttt{name}, \texttt{dim}, \texttt{bounds},
        \texttt{problem\_type}, \texttt{default\_noise\_std}.
        Entry point \texttt{observe(X)} validates inputs via \texttt{prepare\_X} before
        delegating to the abstract \texttt{\_observe\_prepared(X)}.
        Overridable helpers: \texttt{has\_noiseless()},
        \texttt{oracle\_target\_distance(T)}.
\end{itemize}
 
\paragraph{\texttt{common/models.py}.}
\begin{itemize}
  \item \texttt{GPModelConfig}: \texttt{noise\_std} (\texttt{None} $\Rightarrow$
        \texttt{SingleTaskGP}; \texttt{float} $\Rightarrow$ \texttt{FixedNoiseGP}),
        \texttt{normalize\_input}, \texttt{standardize\_output}, \texttt{kernel}.
  \item \texttt{GPModelBuilder}: \texttt{.construct(X, Y)} (no fitting),
        \texttt{.fit(X, Y)} (construct + MLL), \texttt{.predict(model, X)},
        \texttt{.extract\_hyperparameters(model)}.
  \item \texttt{serialize\_state\_dict(model)} and
        \texttt{load\_state\_dict\_into(model, serial)}: JSON-safe round-trip,
        strict key matching.
\end{itemize}
 
\paragraph{\texttt{common/initial\_design.py}.}
Single function: \texttt{sobol\_design(bounds, n, seed)} returns a deterministic
scrambled Sobol grid.
 
\paragraph{\texttt{common/acquisition\_utils.py}.}
Generic wrappers \texttt{optimize\_acquisition} and \texttt{suggest\_next\_point}
around \texttt{botorch.optim.optimize\_acqf}, plus \texttt{clamp\_variance}.
 
\paragraph{\texttt{common/metrics.py}.}
Pure stateless functions: \texttt{target\_distance}, \texttt{target\_best\_so\_far},
\texttt{target\_simple\_regret}, \texttt{target\_cumulative\_regret},
\texttt{hitting\_time}, \texttt{area\_under\_curve}, and the convenience wrapper
\texttt{final\_run\_summary(f\_values, target, oracle\_best, epsilon,
failure\_threshold)} that returns the dictionary consumed by
\texttt{process.py::\_final\_scalars}.
 
\paragraph{\texttt{common/diagnostics.py}.}
\texttt{gaussian\_nlpd}, \texttt{standardized\_residuals},
\texttt{empirical\_coverage}, \texttt{pit\_values},
\texttt{confidence\_bound\_violation\_rate};
combined into \texttt{calibration\_summary(y\_true, mean, std)} and
\texttt{near\_target\_calibration\_summary}.
Joint-posterior diagnostics: \texttt{joint\_gaussian\_nlpd} and
\texttt{posterior\_joint\_calibration\_summary} (Cholesky-stabilised joint NLPD).
The orchestrator \texttt{build\_calibration\_report\_dict} assembles the global,
near-target, joint, and prediction-grid blocks for one snapshot.
Optional KS tests: \texttt{ks\_test\_standard\_normal},
\texttt{ks\_test\_uniform}.

\paragraph{\texttt{common/statistics.py}.}
\texttt{summarize\_sample} (descriptive), \texttt{failure\_rate},
\texttt{wilcoxon\_paired\_comparison} (paired primary test),
\texttt{mann\_whitney\_comparison}, \texttt{rank\_biserial\_correlation},
\texttt{cohens\_d\_right\_better}, \texttt{bootstrap\_ci\_curves} (pointwise),
\texttt{bootstrap\_uniform\_band} (simultaneous), and
\texttt{aggregate\_convergence\_curves} (dispatches on
\texttt{band\_type} $\in\{$\texttt{pointwise}, \texttt{uniform}, \texttt{both}$\}$).
Sign convention: smaller is better; positive effect sizes mean \RIGHT{} beats \WRONG{}.
 
\paragraph{\texttt{common/logging.py}.}
\texttt{RunLogger} owns the run-directory lifecycle:
\texttt{create}, \texttt{save\_config},
\texttt{save\_metadata} (auto-captures git commit, dirty flag, \texttt{sys.argv},
Python version, platform, UTC timestamp),
\texttt{save\_seed\_result},
\texttt{processed\_per\_seed\_path},
\texttt{processed\_aggregate\_path},
\texttt{setup\_file\_logger}.
 
\paragraph{\texttt{common/serialization.py}.}
Pure format conversion: \texttt{tensor\_to\_list}, \texttt{tensors\_to\_lists},
\texttt{save\_json}, \texttt{load\_json}, \texttt{save\_csv}, \texttt{load\_csv},
\texttt{save\_pickle}.
 
\subsection{\texttt{benchmarks/}}
\label{sec:code-benchmarks}
 
\texttt{benchmarks/synthetic.py} defines the abstract \texttt{SyntheticBenchmark} class
extending \texttt{ProblemSpec} with i.i.d.\ Gaussian noise, the \texttt{test\_set}
helper, the \texttt{near\_target\_test\_set} rejection sampler, and a
\texttt{from\_config} entry point (with strict unknown-key checking).
The eight concrete benchmarks (\texttt{SineTrend1D}, \texttt{Forrester1D},
\texttt{Branin2D}, \texttt{Hartmann6D}, \texttt{Linear1D}, \texttt{Exponential1D},
\texttt{GaussianTrap1D}, \texttt{DampedOscillatoryTrap1D}) live in
\texttt{benchmarks/toy.py}.
The module-level \texttt{BENCHMARK\_REGISTRY: dict[str, type[SyntheticBenchmark]]}
and the helper \texttt{build\_benchmark(config)} (both exported from
\texttt{benchmarks/\_\_init\_\_.py}) instantiate a benchmark from a YAML config
(flat or nested layout) or from per-seed metadata.
 
\subsection{\texttt{acquisitions/target\_matching.py}}
\label{sec:code-acq}
 
Two BoTorch acquisitions, both subclasses of \texttt{AnalyticAcquisitionFunction}
with $q = 1$:
 
\begin{lstlisting}[language=Python]
class TargetMatchingUCB(AnalyticAcquisitionFunction):
    """RIGHT. Trained model: GP on raw y_obs ~ f.
       Maximizes: alpha = -|mu_f - T| + sqrt(beta) * sigma_f."""
    def __init__(self, model, target, beta=2.0): ...
 
    @t_batch_mode_transform()
    def forward(self, X):
        posterior = self.model.posterior(X)
        mean  = posterior.mean
        sigma = posterior.variance.clamp_min(1e-10).sqrt()
        return -(mean - self.target).abs() + math.sqrt(self.beta) * sigma
 
 
class WrongTargetMatchingUCB(AnalyticAcquisitionFunction):
    """WRONG. Trained model: GP on g(x) = |y_obs - T|.
       Maximizes: alpha = -mu_g + sqrt(beta) * sigma_g.
       Note: T is baked into training data, not passed here."""
    def __init__(self, model, beta=2.0): ...
 
    @t_batch_mode_transform()
    def forward(self, X):
        posterior = self.model.posterior(X)
        mean  = posterior.mean
        sigma = posterior.variance.clamp_min(1e-10).sqrt()
        return -mean + math.sqrt(self.beta) * sigma
\end{lstlisting}
 
\subsection{\texttt{data/} --- On-Disk JSON Schemas}
 
\texttt{data/raw.py} and \texttt{data/processed.py} contain only dataclasses with
\texttt{to\_dict}/\texttt{from\_dict} round-trip methods.
They depend only on the standard library.
This lets \texttt{process.py} and \texttt{analyze.py} parse on-disk artifacts
symbolically without re-implementing the schema.
 
\subsection{\texttt{experiments/exp\_a\_target\_matching/}}
 
\paragraph{\texttt{run.py}.}
Entry point:
\texttt{python -m experiments.exp\_a\_target\_matching.run
  --config <YAML> [--tag] [--results-root] [--device]}.
\begin{enumerate}
  \item Parse arguments, load YAML, resolve device.
  \item Look up benchmark via \texttt{BENCHMARK\_REGISTRY[name]}.
  \item Create run directory; save config and metadata.
  \item For each seed: \texttt{\_run\_one\_seed} draws the Sobol initial design, obtains
        shared noisy observations, and loops over methods.
        \texttt{\_run\_one\_method} runs the BO loop for \texttt{budget} steps,
        appending one \texttt{StepRecord} and one \texttt{ModelSnapshot} per iteration.
  \item Save seed result via \texttt{logger.save\_seed\_result}.
\end{enumerate}
 
\paragraph{\texttt{process.py}.}
Entry point:
\texttt{python -m experiments.exp\_a\_target\_matching.process
  --run-dir <run\_dir> [--benchmark ...] [--seed ...] [--device cpu|cuda]}.
For each discovered \texttt{raw/<bench>/seed\_NNN.json}:
\begin{enumerate}
  \item Reconstruct the full $N$-step evaluation history.
  \item Compute \texttt{FinalScalars} via \texttt{common.metrics.final\_run\_summary}.
  \item Compute \texttt{ConvergenceCurves} via the four target-distance/regret helpers.
  \item For each retained \texttt{ModelSnapshot}: reconstruct GP via
        \texttt{construct + load\_state\_dict\_into} (no MLL refitting),
        evaluate posterior on $X_{\mathrm{test}}$, compute calibration summaries,
        store \texttt{CalibrationPredictionGrid}.
  \item Wrap in \texttt{ProcessedSeedData} and write per-seed JSON.
\end{enumerate}
 
\paragraph{\texttt{analyze.py}.}
Entry point:
\texttt{python -m experiments.exp\_a\_target\_matching.analyze
  --run-dir <run\_dir> [--benchmark ...] [--statistic mean|median]
  [--confidence FLOAT] [--n-bootstrap INT]}.
For each benchmark, loads all per-seed processed JSON files into
\texttt{ProcessedSeedData} objects and writes the four aggregate JSON files.
 
\subsection{\texttt{scripts/} and \texttt{hpc/narval/}}

\texttt{scripts/} holds convenience wrappers: \texttt{run\_all\_cpu.sh} (run every
config), \texttt{process\_last\_runs.sh} (process + analyze the most recent runs),
and \texttt{smoke\_test\_exp\_a\_workflow.py} (legacy end-to-end smoke wrapper).
\texttt{hpc/narval/} is reserved for SLURM job scripts for the Narval compute
cluster but is currently empty; full-scale runs are launched via the
\texttt{scripts/} wrappers.
 
% ============================================================
\section{Key Experimental Contracts and Invariants}
\label{sec:contracts}
% ============================================================
 
The following contracts are explicitly enforced by the code.
A violation of any contract would invalidate the scientific validity of the pipeline.
 
\begin{contract}[BO trains on $y_{\mathrm{obs}}$ only]
The BO algorithm never sees $f_{\mathrm{noiseless}}$.
\texttt{run.py} feeds \texttt{spec["transform"](Y\_raw, T)} to
\texttt{GPModelBuilder.fit}.
The noiseless value is recorded in \texttt{StepRecord} only for post-hoc analysis.
\end{contract}
 
\begin{contract}[Paired initial design across methods]
\label{contract:paired-ref}
Restates Contract~\ref{contract:paired}.
The \texttt{InitialDesign} is stored at the seed level of the raw JSON, produced once
before any method runs.
\end{contract}
 
\begin{contract}[Calibration uses load-only reconstruction]
Restates Contract~\ref{contract:no-refit}.
\texttt{process.py} rebuilds GPs via \texttt{construct + load\_state\_dict\_into}.
MLL is never refitted. The offline GP is bit-identical to the one used during BO.
\end{contract}
 
\begin{contract}[Raw / processed separation]
Raw files contain trajectory facts only.
\texttt{process.py} reads at most one raw seed file at a time.
\texttt{analyze.py} reads only processed per-seed files, never raw files.
\end{contract}
 
\begin{contract}[Held-out grid is shared and deterministic]
\texttt{benchmark.test\_set(n, seed)} is a pure function of $(\text{benchmark}, n,
\text{seed})$.
The same $X_{\mathrm{test}}$ is used for both methods and across all BO seeds within
a run.
\end{contract}
 
\begin{contract}[Method-aware $y_{\mathrm{train}}$ reconstruction]
\texttt{MethodTrajectory.train\_label} carries the transform recipe
(\texttt{"y\_obs"} vs.\ \texttt{"abs(y\_obs - target)"}).
\texttt{process.py::\_reconstruct\_X\_Y} dispatches on this string and is the single
place that knows the transform.
\end{contract}
 
\begin{contract}[Run provenance]
\texttt{RunLogger.save\_metadata} records UTC timestamp, Python version, platform,
git commit, dirty flag, and \texttt{sys.argv}.
\end{contract}
 
\begin{contract}[Smaller is better, sign convention]
Mann--Whitney is called with \texttt{alternative = "less"}.
Positive $\rho_{\mathrm{rb}}$ and positive Cohen's $d$ both mean \RIGHT{} beats
\WRONG{}.
\end{contract}
 
% ============================================================
\section{Expected Outcome}
\label{sec:expected}
% ============================================================
 
The goal is not merely to show that \RIGHT{} performs better, but to trace the failure
mechanism of \WRONG{} through the full diagnostic chain.
 
\textbf{Expected behavior of \WRONG:}
\begin{itemize}
  \item Distorted posterior behavior near the target level set (kinks in $|f-T|$).
  \item Overconfident uncertainty near kinks; elevated near-target confidence
        violation rate.
  \item Worse calibration, especially in the near-target diagnostic.
  \item Potentially worse final optimization performance.
  \item Sensitivity to the initial Sobol design.
\end{itemize}
 
\textbf{Expected behavior of \RIGHT:}
\begin{itemize}
  \item Smooth GP posterior on the raw function $f$.
  \item Meaningful posterior uncertainty throughout the domain.
  \item Target handling at the acquisition level, without distorting the surrogate.
  \item More graceful degradation across initializations and noise realizations.
\end{itemize}
 
% ============================================================
\section{Future Extensions and Planned Additions}
\label{sec:future}
% ============================================================
 
The following items are explicitly planned but not yet implemented.
They are listed here to prevent confusion between the current state and the roadmap.
 
Several items previously listed here are now \emph{implemented} and described in
the body: the paired Wilcoxon signed-rank test
(\S\,\ref{sec:stats}), the simultaneous (uniform) bootstrap bands
(\S\,\ref{sec:aggregate}), and the additional diagnostic benchmarks
(\S\,\ref{sec:expA}). The remaining roadmap items are:

\begin{description}

  \item[\texttt{make\_figures.py}.]
  A dedicated script for generating publication-quality figures
  (Figures A1--A4) from the processed and aggregate outputs.
  Currently, all figures are produced in the analysis notebooks
  (\texttt{notebooks/Exp\_A\_analysis\_*.ipynb}).

  \item[Target-aware Expected Improvement (target-EI).]
  A more refined SoC-compliant acquisition computing
  $\mathrm{EI}_T(x) = \mathbb{E}_t[\max\{0, \tau_t - |f(x)-T|\}]$
  under the posterior on the raw $f$.
  A Monte Carlo implementation sketch exists in the older design document.
  To be added to \texttt{acquisitions/target\_matching.py} once the main pipeline is
  stable.

  \item[Full-scale paper seeds.]
  The current configurations use 100 seeds; the final paper configuration may use
  a larger number (target TBD).

  \item[Full-scale benchmark sweep.]
  All eight benchmarks have configs in \texttt{configs/}; the remaining work is to
  run, process, and analyze the full sweep and consolidate the cross-benchmark
  tables.

  \item[Experiments B and C.]
  The codebase is designed to support multi-objective SoC (Experiment B) and
  prior-knowledge SoC (Experiment C).
  \texttt{acquisitions/prior\_guided.py}, \texttt{acquisitions/multi\_objective.py},
  and the corresponding experiment directories are not yet implemented.

  \item[Real-data problems.]
  \texttt{common/real\_data.py} is planned (not yet created).
  Real-data experiments are future work after Experiment A is validated.

\end{description}
 
% ============================================================
% Extension cookbook (appendix-style)
% ============================================================
\section*{Appendix: Extension Cookbook}
\addcontentsline{toc}{section}{Appendix: Extension Cookbook}
\label{sec:cookbook}
 
\begin{description}
 
  \item[New benchmark.]
  Add a class to \texttt{benchmarks/toy.py} inheriting \texttt{SyntheticBenchmark};
  implement \texttt{.f(X)} and optionally \texttt{.oracle\_target\_distance(T)};
  register in \texttt{BENCHMARK\_REGISTRY}.
 
  \item[New method.]
  Add the acquisition class to \texttt{acquisitions/};
  add a spec entry to \texttt{\_METHOD\_SPECS} in \texttt{run.py} with
  \texttt{train\_label}, \texttt{transform}, \texttt{acq\_builder},
  \texttt{use\_fixed\_noise}.
  Extend \texttt{process.py::\_reconstruct\_X\_Y} if \texttt{train\_label} is new.
  Use the \texttt{right\_*}/\texttt{wrong\_*} prefix so \texttt{analyze.py} pairs it
  automatically.
 
  \item[New final scalar metric.]
  Extend \texttt{common.metrics.final\_run\_summary};
  add the field to \texttt{data.processed.FinalScalars};
  populate in \texttt{process.py::\_final\_scalars};
  add the key to \texttt{analyze.py::\_FINAL\_SCALAR\_KEYS}.
 
  \item[New calibration diagnostic.]
  Add the function to \texttt{common.diagnostics} and include it in the
  \texttt{calibration\_summary} dictionary;
  add the key to \texttt{\_CALIB\_TEST\_GRID\_KEYS} in \texttt{analyze.py}.
 
  \item[New statistical test.]
  Add a pure function to \texttt{common.statistics};
  call it from \texttt{analyze.py::\_build\_statistical\_tests}.
 
  \item[New configuration field.]
  Add the key to the YAML;
  read it via \texttt{config.get(...)} in the appropriate stage;
  store essential provenance in \texttt{metadata.json}.
 
\end{description}
 
\end{document}